% Generated draft: requires mathematical and source review.
\documentclass[11pt]{article}
\usepackage[margin=25mm]{geometry}
\usepackage{fontspec}
\setmainfont{Latin Modern Roman}
\usepackage{amsmath,amssymb,mathtools,unicode-math}
\setmathfont{Latin Modern Math}
\usepackage{enumitem,needspace,xurl,hyperref}
\setlength{\parindent}{0pt}
\setlength{\parskip}{5pt}
\setlength{\emergencystretch}{3em}
\newcommand{\N}{\mathbb{N}}
\newcommand{\Z}{\mathbb{Z}}
\newcommand{\Q}{\mathbb{Q}}
\newcommand{\R}{\mathbb{R}}
\newcommand{\C}{\mathbb{C}}
\newcommand{\F}{\mathbb{F}}
\newcommand{\A}{\mathbb{A}}
\newcommand{\PP}{\mathbb{P}}
\newcommand{\E}{\mathbb{E}}
\newcommand{\eps}{\varepsilon}
\newcommand{\dd}{\,\mathrm{d}}
\begin{document}
\section*{UnsolvedMath ID: 20001271}
\textbf{Source problem code:} catalogue entry
\subsection{Definitions and mathematical statement}
In the usual nontrivial projection setting, let $K\subset\R^n$ be a convex body (compact, convex, with nonempty interior) and let $2\le k<n$. Let $G(n,k)$ be the Grassmannian of $k$-dimensional linear subspaces and $K|S$ the orthogonal projection. The question asks whether
\[
\forall S,T\in G(n,k)\ \exists\text{ an affine bijection }\phi:S\to T
\text{ with }\phi(K|S)=K|T
\]
forces $K$ to be an ellipsoid, a translate of an invertible linear image of the Euclidean unit ball. No central symmetry is assumed. The range matters: for $k=1$ all projections are intervals, and for $k=n$ the hypothesis is vacuous, so those literal readings would not characterize ellipsoids. The assigned excerpt does not reproduce all conventions from its preceding problem.

The attached aid reports affirmative literature results for proper projection dimensions $k\equiv0,1,2\pmod4$, with possible exception $k=133$, and records $k\equiv3\pmod4$ and $k=133$ as unresolved cases in its assessment. Exact hypotheses and references are not supplied here.

Its geometric endgame is the following more general criterion. Let $\mathcal F$ be a family of linear subspaces that covers every $2$-plane by containment:
\[
\forall L\in G(n,2)\ \exists S\in\mathcal F\text{ with }L\subseteq S.
\]
If $K|S$ is an ellipsoid in $S$ for every $S\in\mathcal F$, then $K$ is an ellipsoid, without assumed central symmetry. The aid describes compatibility of projected centers through the odd part of the support function $h_K(u)=\max_{x\in K}\langle x,u\rangle$: its restriction to each $2$-plane is linear, forcing one global center. The squared centered support function then satisfies the parallelogram law and comes from a positive quadratic form.

Consequently, under the original pairwise affine-equivalence hypothesis with $k\ge2$, a single ellipsoidal $k$-projection makes every $k$-projection ellipsoidal, and the criterion yields an ellipsoidal $K$. The existence of such a projection is an extra intermediate condition, not something proved in the remaining cases by this endgame alone.
\subsection{Short English statement}
Do mutually affinely equivalent proper projections force a convex body to be an ellipsoid? The supplied aid gives a general two-plane-covering endgame and reported dimension-dependent results, while leaving the remaining existence-of-an-ellipsoidal-projection step unresolved.
\subsection{Sources}
\begin{itemize}
\item Local UnsolvedMath registry, permanent ID $20001271$, source code catalogue entry; canonical \texttt{aim-convex-geometry-notes.json}, \texttt{notes[2]}.
\item American Institute of Mathematics, 2013 workshop ``Sections of convex bodies,'' problem 2: \url{https://aimath.org/pastworkshops/sectionsconvexproblems.pdf}. The original refers to problem 1 and to [G, Problem 3.3 and Note 3.2], without supplying that bibliographic expansion or the preceding conventions. No PDF inspection is claimed here.
\item Attached research attempt 1: partial result, run and registry status partially solved, rigor proved, novelty confidence low, classification PARTIAL-PROGRESS. Its literature assessment and two-covering criterion are machine generated and not peer reviewed; independently verify before citation. No authoritative exact references for the congruence cases are supplied.
\end{itemize}
\end{document}
